\section{Clean children and independent repetition}
\label{sec:clean}

A normalized product form of bounded rank determines a short odd list
of answers at each source leaf.  When the decoded forms agree after
projection, their lists satisfy a projection constraint.  To apply
parallel repetition to these lists, we must recover independent source
edges and local answer rules despite their sampled advice.

We do this by finding projected children whose contributions to every
sampled function are zero.  The probability of this event is exactly
independent of their source questions.  Revealing the remaining data
therefore leaves independent source edges at designated leaves, and
each player can reconstruct its original decoder input from its own
questions there.  Repetition then bounds the low-rank agreement from
Proposition~\ref{prop:low-rank}.

Fix levels $i<j$ and use the modified experiment at $D_i$ throughout
this section.  Write
\[
 H_i=H_{D_i}(E),\qquad H_i'=H_{D_i}(O),\qquad S=M_{D_i}.
\]
We identify $H_i'$ with its pullback subspace of $H_i$.  The integer
$s\ge1$ is a full-rank cutoff as in
Proposition~\ref{prop:low-rank}.  The source game, denoted by
$G_{\mathrm{src}}$, is a fixed $t$-fold power of the clause--variable
game.  For a source question $u$ on the left and $v$ on the right,
write $\Lambda_u$ and $\Gamma_v$ for their nonempty legal answer sets.
Thus $\Gamma_v=\F^t$, with a separate coordinate for each slot,
including slots with repeated variable IDs.  Let $\mu$ be the
distribution on source edge records $e=(u,v,\pi_e)$.

\subsection{Odd lists and leaf indicators}

For a finite nonempty set $\Lambda$, define
\[
 \mathcal O_s(\Lambda)
 =\{L\subseteq\Lambda:1\le |L|\le s, |L|\text{ odd}\}.
\]
Given a map $\pi:\Lambda\to\Gamma$, define its parity pushforward by
\begin{equation}
 \pi_{\mathrm{odd}}(L)
 =\{b\in\Gamma:|L\cap\pi^{-1}(b)|\text{ is odd}\}.
 \label{eq:clean-parity-map}
\end{equation}
The same parity pushforward describes restriction of Fourier characters
in low-degree long-code analysis
\cite[Section~2.3, Equation~(1), ECCC version]{DinurGuruswami2015}.
The odd-list game $G_{\mathrm{odd},s}$ has the same questions and edge
distribution as $G_{\mathrm{src}}$, legal answers
$\mathcal O_s(\Lambda_u)$ and $\mathcal O_s(\Gamma_v)$, and edge map
$(\pi_e)_{\mathrm{odd}}$.

\begin{lemma}[Odd-list projection and extraction]
\label{lem:parity-game}
Every map in Equation~\eqref{eq:clean-parity-map} sends
$\mathcal O_s(\Lambda)$ into $\mathcal O_s(\Gamma)$.  In particular,
$G_{\mathrm{odd},s}$ is a total projection game, and
\begin{equation}
 \val(G_{\mathrm{odd},s})\le s^2\val(G_{\mathrm{src}}).
 \label{eq:clean-list-value}
\end{equation}
Suppose a bilinear form $F$ on $H_i$ has the form
\[
 F(x,y)=z(xy),\qquad z\in(H_i^2)^*,\qquad z(1)=1,
 \qquad \rank F\le s.
\]
For every leaf $\lambda$ below $D_i$, with local answer set
$\Lambda_\lambda$, its diagonal on the leaf indicators specifies a
member of $\mathcal O_s(\Lambda_\lambda)$.  If that leaf is projected
by $\pi_e$ and a form $F'$ on $H_i'$ equals the restriction of $F$,
then its corresponding diagonal list is the parity pushforward of
the left list.
\end{lemma}

\begin{proof}
For $L\in\mathcal O_s(\Lambda)$, summing all preimage multiplicities
modulo two gives
\[
 |\pi_{\mathrm{odd}}(L)|\equiv |L|\equiv1\pmod2.
\]
Each member of $\pi_{\mathrm{odd}}(L)$ has at least one preimage in
$L$, and these preimage sets are disjoint.  Consequently its size is
at most $|L|$, proving totality.

Take any deterministic labeling of the odd-list game.  Each player
independently chooses a uniformly random member of the list it would
return.  On an accepted list edge, the two lists contain an accepting
ordinary answer pair: any element of the nonempty right list has an
odd, hence positive, number of preimages in the left list.  The
independent choices select some accepting pair with probability at
least $1/(|L||R|)\ge s^{-2}$.  This is a legal randomized source
strategy; averaging deterministic answer tables bounds its success
by $\val(G_{\mathrm{src}})$.  Maximizing over the list labeling proves
Equation~\eqref{eq:clean-list-value}.

For the extraction statement, let $e_a$ be the indicator that the
answer at leaf $\lambda$ is $a\in\Lambda_\lambda$, extended by
ignoring every other leaf.  By Definition~\ref{def:spaces},
$e_a\in H_i$.  Pointwise multiplication gives
\[
 e_ae_b=0\quad(a\ne b),\qquad e_a^2=e_a,
 \qquad \sum_{a\in\Lambda_\lambda}e_a=1.
\]
Thus $(F(e_a,e_b))_{a,b}$ is diagonal, with diagonal
$d_a=z(e_a)$.  Its rank is the integer $|\{a:d_a=1\}|$ and is at
most $\rank F\le s$.  Moreover,
$\sum_a d_a=z(1)=1$ in $\F$, so this support is odd and nonempty.

For a projected answer $b$, the pullback of its indicator is
\[
 e_b'=\sum_{a:\pi_e(a)=b}e_a.
\]
Agreement of the forms and the diagonal identity imply
\begin{equation}
 F'(e_b',e_b')
 =\sum_{a:\pi_e(a)=b}d_a.
 \label{eq:clean-indicator-pushforward}
\end{equation}
The support on the right is exactly the parity pushforward in
Equation~\eqref{eq:clean-parity-map}, completing the proof.
\end{proof}

Low-rank symmetric matrices also yield short assignment lists in the
quadratic-code analysis of Khot and Saket, through rank-one decompositions
\cite[Lemma~2.1 and Section~7.2, ECCC version]{KhotSaket2017}.
Here the orthogonal leaf indicators give the lists directly as diagonal
supports of legal answers.

We use a total local extraction rule also when a decoder fails to
return such a form.  Order the legal answers to every source question
once and for all.  At each designated leaf, a player reads the
diagonal support of its own output form.  If that support is odd and
has size at most $s$, it returns the support.  If the output is not a
form or this test fails, it returns the singleton consisting of the
first legal answer to its own question.  The test and default use
only that player's input and output.  In particular, the right
player need not determine the rank or validity of the left form.
These rules produce legal lists on every input, and Lemma~\ref{lem:parity-game}
identifies their outputs whenever normalized low-rank agreement occurs.

\subsection{An exact law for vanishing advice}

Zero advice was already used to hide a source coordinate and recover
soundness under repetition by Khot, Minzer and Safra
\cite[Sections~3.3--3.4, especially Remark~3.2]{KhotMinzerSafra2017}.
Here we zero all sampled contributions from selected children and prove
that the resulting conditional source law is exactly a product law.

Choose the first leaf of each child of $D_i$ as its designated leaf.
This choice depends only on the ordered tree.  We will express every
uniform cut draw as a sum of independent child contributions.  A projected
child will be clean when all its contributions and all its internal arrays
are zero.  Its arguments then disappear from the sampled functions, while
the zero probability can be computed on that child's projected domain.

The family of contributions must cover every input used by the decoders.
Stop every recursive scalar call used by an ancestor array when it first
reaches $D_i$, including the calls for every bucket of $M=M_{D_j}$.
Fix this family before the row map $A$ selects visible buckets, and add
the $\ell_i$ independent rows of the uniform matrix $S$.
By Equation~\eqref{eq:direct-cut-count}, its size is
\begin{equation}
 T_i^{\mathrm{cl}}=\widehat T_i=\ell_i+
 \sum_{h=i+1}^{d}(2^{\ell_h}-1)
                   \prod_{p=i+1}^{h}(2R_p).
 \label{eq:clean-call-count}
\end{equation}
Indeed, an array at level $h$ has $2^{\ell_h}-1$ scalar buckets, and
each scalar call at level $p>i$ makes $2R_p$ calls in its path child.
All other contributions lie outside $D_i$, and distinct calls use fresh
randomness.  If a recipe uses fewer calls, pad it with independent unused
calls to reach $T_i^{\mathrm{cl}}$.  Thus the call count is fixed before any
visibility decision.  The notation $T_i^{\mathrm{cl}}$ distinguishes this
family from those enlarged by extra rows in comparison experiments.

Let $\Xi$ record the questions outside $D_i$, the path through $D_i$,
the exterior ingredients of the stopped expressions, and all other arrays
supported outside $D_i$.  Include the independent choices $A$ and $a_i$
after fixing the call family.  Record each exterior function on the domain
of its supporting subtree, together with its location in the expression.
This representation requires no endpoint questions inside $D_i$; each
player will later extend it to its own product domain by ignoring the
other arguments.  Assembled ancestor values are still functions of the
pending cut calls and are not exposed in $\Xi$.
Consequently $\Xi$ depends only on exterior source records and independent
exterior randomness, so the source samples inside $D_i$ retain their
independent laws conditional on $\Xi$.

It is convenient to sample a complete source edge at every leaf,
even when its child will remain unprojected.  This merely augments
the modified experiment: sample the left question and then a source
projection conditional on it, using that projection only if the
child is selected.  For each child $C$ of $D_i$, independently let
$P_C$ be its projection flag, with
$\Pr(P_C=1)=\beta_i=n_i^{-2/3}$.  Write $\boldsymbol e_C$ for the
tuple of source edge records at its leaves and $m_i$ for the number
of those leaves.  Thus $\boldsymbol e_C$ has distribution
$\mu^{\otimes m_i}$.

Generate cut call $q\in[T_i^{\mathrm{cl}}]$ by taking, independently for all
$q,C$, a uniform latent summand
\[
 U_{q,C}\in H_C(O)^2,
 \qquad U_q=\sum_{C\text{ child of }D_i}U_{q,C}.
\]
The sum is uniform in $H_i'$ because the addition map from the
product of the child square spaces onto $H_i'$ is a surjective linear
map.  Its fibers all have the same size.  The resulting calls are
independent, and are independent of the uniform arrays at every
nonleaf proper descendant of $D_i$.  The first $\ell_i$ calls can
therefore serve as the rows of $S$, with the rest substituted in the
stopped ancestor expressions.  This is exactly the modified sampling
law on the visible arrays.  No direct-sum assumption on the child
spaces is involved.

Call $C$ \emph{clean} if $P_C=1$, every $U_{q,C}$ is zero, and every
sampled array at a nonleaf node in the subtree rooted at $C$ is zero.
Write $\mathcal C$ for the set of clean children.

\begin{lemma}[Exact clean probability and conditional source law]
\label{lem:clean-law}
Conditional on any exterior datum $\Xi=\xi$ of positive probability,
the clean indicators are independent.  In the padded construction
above each has probability $\beta_i p_{0,i}$, where
\begin{equation}
 p_{0,i}=
 2^{-T_i^{\mathrm{cl}}\dim(H_C'^2)
       -\displaystyle\sum_{\substack{D\preceq C\\D\text{ nonleaf}}}
           \ell_{\operatorname{level}(D)}\dim H_D'}>0.
 \label{eq:clean-zero-probability}
\end{equation}
Here $D\preceq C$ includes $D=C$, and $H_D'$ is the function space
on a fully projected subtree of the indicated shape.  The exponent
has the same value for every child and every choice of source
questions and projection positions.  It depends only on the source
power, the fixed row and product counts, and branching below level
$i$; it is independent of the instance, $n_i$, and all higher
branching numbers.

After fixing $\mathcal C$, expose all source questions, projection
choices and latent data at nonclean children, and all source edge
records at nondesignated leaves of clean children.  Denote this
additional exposure together with $\Xi$ and $\mathcal C$ by
$\mathcal D_{\mathrm{pub}}$.  Conditional on every value of $\mathcal D_{\mathrm{pub}}$ of positive
probability, the source edge records at the designated clean leaves
are independent and each has distribution $\mu$.
\end{lemma}

\begin{proof}
Fix $\Xi=\xi$ and a child edge tuple $\boldsymbol e_C$.  Conditional
on $P_C=1$, the $T_i^{\mathrm{cl}}$ latent summands in that child are independent
uniform points of $H_C'^2$.  Their probability of all being zero is
$2^{-T_i^{\mathrm{cl}}\dim(H_C'^2)}$.  An independent uniform
$\ell_p$-row array in $H_D'$ is zero with probability
$2^{-\ell_p\dim H_D'}$.  Multiplication over all such descendant
arrays proves Equation~\eqref{eq:clean-zero-probability}.
This calculation sets each latent summand to zero separately.
The possibility that nonzero summands could cancel, including shared
constants, has no effect on this event or on its probability.

We next prove that the displayed probability is exactly independent
of $\boldsymbol e_C$.  Each fully projected leaf has answer domain
$\F^t$.  Given two projected subtrees of the same ordered shape,
identify these domains coordinatewise, retaining a separate factor
at every leaf and every source-power slot.  Pullback by this product
bijection identifies their full leaf function spaces.  It preserves
sums and pointwise products, so induction through the tree identifies
every corresponding $H_D'$ and $H_D'^2$.  In particular their
dimensions agree.  Neither the question IDs nor the selected
projection positions enter these identifications.  Even a right
question $(v,v)$ has the full answer domain $\F^2$ in its two slots;
it does not identify their answer coordinates.

Consequently, for every possible edge tuple $\boldsymbol e$,
\begin{equation}
 \Pr(\boldsymbol e_C=\boldsymbol e, C\in\mathcal C\mid\Xi=\xi)
 =\mu^{\otimes m_i}(\boldsymbol e)\,\beta_i p_{0,i}.
 \label{eq:clean-factorization}
\end{equation}
The latent child blocks, edge tuples and flags are independent
across children conditional on $\Xi$.  Thus the clean indicators are
independent, and the constant factor in
Equation~\eqref{eq:clean-factorization} cancels on conditioning that
a child is clean.  The entire edge tuple in each clean child retains
its original product law.  Positivity of a question-dependent zero
probability would not suffice for this cancellation; its constancy
has just been proved.

Fixing the whole clean mask still leaves a product across children.
Exposing full data of nonclean children therefore imposes no
additional condition on any clean edge tuple.  Within each clean
child, fixing the nondesignated edge records in its product tuple
leaves the designated edge with law $\mu$.  This proves the asserted
law conditional on $\mathcal D_{\mathrm{pub}}$.  Natural repetitions of an ID in
different samples are permitted throughout; this exposure neither
requires nor forbids a collision with an exposed ID.

Finally, the child shapes involve only $n_1,\ldots,n_{i-1}$, and
Equation~\eqref{eq:clean-call-count} involves row and product counts
but no branching numbers.  This proves the stated parameter
dependence.  If $i=1$, the children are leaves and the descendant-array
sum is empty, with value zero, so the same argument applies.
\end{proof}

\subsection{Reconstructing the original local strategies}

The preceding lemma leaves independent source questions, but repetition
also requires local answer rules on those questions.  Here the input
restriction is substantive: the left decoder in
Proposition~\ref{prop:upper} never receives $a_i$, while the right
decoder receives $S$ only modulo $a_i$.  We now reconstruct precisely
those inputs after the final exposure.

\begin{lemma}[Local reconstruction on clean children]
\label{lem:reconstruction}
Fix $\mathcal D_{\mathrm{pub}}=d_0$ of positive probability as in
Lemma~\ref{lem:clean-law}.  From $d_0$ and its own designated source
questions in $\mathcal C$, each player can reconstruct its input to
the original decoder of Proposition~\ref{prop:upper}.  Composing
these decoders with the total diagonal extraction rule defines local
strategies for $G_{\mathrm{odd},s}^{\otimes |\mathcal C|}$.
Whenever their original forms agree on $H_i'$ and the left form is
normalized, of product form, and of full rank at most $s$, these
strategies accept every designated clean edge.
\end{lemma}

\begin{proof}
First reconstruct the two question tuples.  The left player combines
its designated clause-tuple questions with all the exposed questions
to obtain $E$.  It does not need the selected positions or variable
IDs at its designated clean leaves.  The right player combines its
designated variable-tuple questions with the exposure to obtain $O$:
every clean child is projected, and the entire question and projection
data at every nonclean child are exposed.  The right player does not
need the missing clause-tuple questions.  Repeated IDs specify
separate slots in both reconstructions.

We next reconstruct the functions appearing in the advice.  At a
nonclean child the exposure supplies its actual latent summands and
descendant arrays, including the projections needed to pull them
back on the left.  At a clean child substitute a zero symbol for
each latent summand and descendant row.  Such a symbol specifies
only its call, row and tree location.  The left interprets it as the
zero function on its own original domain; the right interprets it
as the zero function on its own projected domain.  It carries no
missing endpoint ID, selected position or partner-domain metadata.
In either interpretation it equals the corresponding actual
function, because zero pulls back to zero under every projection.

Sum these child contributions to obtain every cut call, including
all rows of $S$.  Then evaluate the stopped expressions using the
exposed exterior ingredients.  This yields all ancestor arrays and
all their individual buckets, hence every nonhidden bucket required
by the decoders.  This is a reconstruction from the chosen sampled
decomposition, not an attempt to recover a unique decomposition from
the resulting function.  Overlaps of constants between child
spaces, or other nonuniqueness of a sum representation, do not
alter the resulting sums.

To check that no missing domain information enters these expressions,
consider the full algebra of functions on the product of all leaf
domains that ignore the arguments in clean children.  It contains
constants and is closed under sums and pointwise products.  Every
reconstructed cut call belongs to this algebra: its clean summands
are zero and its other summands are supported on nonclean children.
Every exterior ingredient also belongs to it, since its support is
outside $D_i$.  Induction over each stopped expression therefore
shows that every reconstructed ancestor function ignores all clean
arguments.  Descendant arrays in clean children are identically
zero.  Thus products at higher levels cannot introduce dependence
on the unavailable endpoint questions.

Each player now has its own domains and all the functions supplied
to its decoder.  It can form their reduced-support partition keys,
consult the fixed labeling, and restore the evaluations of the
selected parts in its own domain, using the conventions of
Section~\ref{sec:construction}.  For each reconstructed sampled joint function, every clean slot is
unused, so canonicalization discards that slot and its domain metadata.  On nonclean projected
slots, the exposed projection supplies the necessary pullback; on
designated clean slots, the functions just proved independent of
those arguments require none.  A decoder can also evaluate any
hypothetical matrices called for by its algorithm: it constructs
them on its own now-known domain.  Such additional matrices need
not ignore the designated arguments, since the player already knows
its own questions there.

Although this final public record may determine full $S$, the
simulated right decoder is called with its original input
$S\bmod a_i$.  Similarly the left decoder is called without $a_i$,
even though the final public record may contain that direction.
Neither decoder receives the reconstructed hidden $M$ or additional
buckets that are absent from its specified input.  Extra public
information in the simulation is simply omitted from the relevant
arguments.  All information used to form the original arguments is
a function of $d_0$ and the player's own designated questions.

The simulation runs the same fixed algorithms on these reconstructed
inputs.  The left uses its fixed usefulness marks and deterministic witness
rule.  The right computes the heavy lists in
Equation~\eqref{eq:upper-heavy-definition} using the original law and response
function determined by its own input in Section~\ref{sec:right-hidden-law},
then applies the same selection and contraction rules.
Conditioning on $\mathcal D_{\mathrm{pub}}$ may change the physical distribution of the
sampled hidden buckets, but it does not change the reference law prescribed
by this algorithm.  The local simulation therefore requires no small-bias
assertion under the clean conditioning.

Finally, run the total diagonal extraction rule at the designated
clean leaves.  Each player has exactly its own source question there,
so this is a legal local answer rule on the repeated game.  The rules
are defined even on inputs on which the original decoder chooses a
default, produces a high-rank form, or fails the diagonal test.  On
normalized low-rank agreement, Lemma~\ref{lem:parity-game} says that
the left supports are valid odd lists and their right supports are
their parity pushforwards.  Both diagonal tests then pass, so every
designated clean edge accepts.  This is an event inclusion, with no
conditioning on the forms' validity, their agreement, or their rank.
\end{proof}

\subsection{The vanishing upper bound}

We can now bound low-rank agreement in the original modified law.
All public exposures used for this purpose have already been
specified, and none includes a useful-advice event, a slice event,
a decoder output or a successful character match.

\begin{proposition}[Clean repetition bound]
\label{prop:clean-bound}
Assume
$\val(G_{\mathrm{src}})<1/(2s^2)$.  Let
$\mathcal L_{i,j,s}$ be the event in the modified experiment that the
two original decoder outputs agree on $H_i'$, while the left output
is a normalized product form on $H_i$ of full rank at most $s$.
Then, for the clean probability in
Equation~\eqref{eq:clean-zero-probability},
\begin{equation}
 \Pr(\mathcal L_{i,j,s})
 \le \left(1-\frac{\beta_i p_{0,i}}{64}\right)^{n_i}
 \le \exp\!\left(-\frac{p_{0,i}}{64}n_i^{1/3}\right).
 \label{eq:clean-final-bound}
\end{equation}
The bound holds for randomized decoders and is uniform over the
source instance and over the fixed labeling.
\end{proposition}

\begin{proof}
Lemma~\ref{lem:parity-game} gives
$\val(G_{\mathrm{odd},s})<1/2$.  Fix a public record
$\mathcal D_{\mathrm{pub}}=d_0$, and put $k=|\mathcal C|$.  By
Lemma~\ref{lem:clean-law}, its designated edge records have exactly
the independent product distribution of $k$ source samples.  By
Lemma~\ref{lem:reconstruction}, the original decoders and the total
extraction rules give local strategies for the corresponding odd-list
game.  Theorem~\ref{thm:repetition}, with gap $1/2$, bounds their
probability of accepting all $k$ edges by $(63/64)^k$.
For $k=0$ the bound is interpreted as one.  Private or shared random
tapes independent of the questions cause no difficulty: fixing
complete tapes gives deterministic local functions, to each of which
the same bound applies, and then one averages over those tapes.

The event inclusion in Lemma~\ref{lem:reconstruction} therefore gives
\[
 \Pr(\mathcal L_{i,j,s}\mid\mathcal D_{\mathrm{pub}}=d_0)
 \le (63/64)^{|\mathcal C|}.
\]
This inequality bounds the favorable agreement event; that event
has never been used to define the conditional source distribution
or the local strategies.  Average over $\mathcal D_{\mathrm{pub}}$, and then use
independence of the clean indicators conditional on $\Xi$.  In the
padded construction this yields the exact generating-function identity
\[
 \E\bigl[(63/64)^{|\mathcal C|}\mid\Xi\bigr]
 =\prod_{C\text{ child of }D_i}
       \left(1-\beta_i p_{0,i}
                   +\frac{63}{64}\beta_i p_{0,i}\right)
 =\left(1-\frac{\beta_i p_{0,i}}{64}\right)^{n_i}.
\]
The second bound in Equation~\eqref{eq:clean-final-bound} follows from
$1-x\le e^{-x}$ and $n_i\beta_i=n_i^{1/3}$.
\end{proof}

\begin{remark}[Exterior-dependent call families]
The padded family above gives one constant clean probability.  The argument
also permits independent exterior recipe choices with
$T_i^{\mathrm{cl}}(\xi)\le T_i^{\mathrm{cl}}$ calls.  Substituting this
count in Equation~\eqref{eq:clean-zero-probability} gives an exact
question-independent probability $p_{0,i}(\xi)\ge p_{0,i}$.
If the recipe differs between children, denote these probabilities by
$p_{0,C}(\xi)$.  Conditional on $\Xi=\xi$, the clean indicators remain
independent, and every $p_{0,C}(\xi)$ is at least $p_{0,i}$.
The exact factors, rather than this lower bound, cancel in
Equation~\eqref{eq:clean-factorization}; hence the conditional source law
and the local reconstruction are unchanged.

The same generating-function calculation then gives
\begin{equation}
 \Pr(\mathcal L_{i,j,s})
 \le \E_\Xi\prod_{C\text{ child of }D_i}
       \left(1-\frac{\beta_i p_{0,C}(\Xi)}{64}\right)
 \le \exp\!\left(-\frac{p_{0,i}}{64}n_i^{1/3}\right).
 \label{eq:clean-exterior-average}
\end{equation}
Only the second inequality uses the uniform positive lower bound.
\end{remark}

For the decoders supplied by Proposition~\ref{prop:upper},
Proposition~\ref{prop:low-rank} gives
$\Pr(\mathcal L_{i,j,s})\ge\gamma_j/2$.  Once the source power and
the lower branching numbers have been fixed,
$p_{0,i}>0$ is already fixed independently of $n_i$ and of higher
branching.  Taking, for example,
\[
 n_i^{1/3}>\frac{64}{p_{0,i}}\log\frac{4}{\gamma_j}
\]
makes Equation~\eqref{eq:clean-final-bound} smaller than
$\gamma_j/4$, a contradiction.  There are only finitely many choices
of $j>i$, so one branching number at level $i$ accommodates all of
them.  Enlarging a higher branching number adds only exterior data
and changes the deterministic stopped expressions; it changes
neither the number of calls in Equation~\eqref{eq:clean-call-count}
nor the projected child dimensions in
Equation~\eqref{eq:clean-zero-probability}.  This is the uniform
contradiction needed for the parameter selection in
Section~\ref{sec:assembly}.

